\documentclass[10pt,reqno]{amsart}

\usepackage[T1]{fontenc}
\usepackage[utf8]{inputenc}
\usepackage{newtxtext,newtxmath}
\usepackage{microtype}
\usepackage{mathtools}
\usepackage{enumitem}
\usepackage[nameinlink,noabbrev]{cleveref}
\usepackage[a4paper,margin=28mm,headheight=15pt]{geometry}

\setlength{\parskip}{2pt}
\setlist[enumerate]{leftmargin=7mm,itemsep=2pt,topsep=3pt}
\emergencystretch=2em

\newtheorem{theorem}{Theorem}[section]
\newtheorem{proposition}[theorem]{Proposition}
\newtheorem{lemma}[theorem]{Lemma}
\newtheorem{corollary}[theorem]{Corollary}
\theoremstyle{definition}
\newtheorem{definition}[theorem]{Definition}
\newtheorem{remark}[theorem]{Remark}

\title[A Cayley graph on \(\mathbb Z_{15}\) with core \(C_5\)]
{A Cayley Graph on \(\mathbb Z_{15}\) Whose Core Is \(C_5\)}

\author{Argus Open Research Project}
\address{Argus Open Research Project}
\date{August 30, 2026}

\subjclass[2020]{Primary 05C60; Secondary 05C25, 20K01}
\keywords{Cayley graph, graph core, graph homomorphism, retraction}

\begin{document}

\begin{abstract}
We give an explicit connected Cayley graph on the cyclic group
\(\mathbb Z_{15}\) whose core is the five-cycle.  Let
\[
 X=\operatorname{Cay}\bigl(\mathbb Z_{15},
 \{1,4,6,9,11,14\}\bigr).
\]
Reduction modulo \(5\) maps \(X\) homomorphically onto \(C_5\), and the map
\(i\mapsto6i\) gives a section whose image is an induced five-cycle.  The
resulting composite is a retraction.  Every endomorphism of \(C_5\) is an
automorphism, so \(C_5\) is the core of \(X\).  Since a Cayley graph on
\((\mathbb Z_{15})^m\) has \(15^m\) vertices, this core is not a Cayley graph
on any power of \(\mathbb Z_{15}\).  The example refutes the literal
arbitrary-abelian formulation of a problem on cores of Cayley graphs, but
does not address prime, elementary-abelian, or cubelike variants.
\end{abstract}

\maketitle

\section{The graph}

Let
\[
 G=\mathbb Z_{15},\qquad
 \pi:G\longrightarrow\mathbb Z_5,\qquad \pi(x)=x\bmod5,
\]
and put
\[
 S=\pi^{-1}(\{1,-1\})
 =\{1,4,6,9,11,14\}.
\]
Define the simple undirected Cayley graph
\[
 X=\operatorname{Cay}(G,S)
\]
by declaring \(x\) adjacent to \(y\) precisely when
\(y-x\pmod {15}\) lies in \(S\).

The set \(S\) does not contain zero and satisfies \(S=-S\), so \(X\) is a
simple undirected graph.  It is \(6\)-regular and has
\[
 \frac{15\cdot6}{2}=45
\]
edges.  Since \(1\in S\), the edges \(x\sim x+1\) contain a spanning
fifteen-cycle, and hence \(X\) is connected.

\section{A retraction onto an induced five-cycle}

Write
\[
 C_5=\operatorname{Cay}(\mathbb Z_5,\{1,-1\})
\]
and define
\[
 \sigma:\mathbb Z_5\longrightarrow\mathbb Z_{15},
 \qquad \sigma(i)=6i\bmod15.
\]
Since \(6\equiv1\pmod5\), one has
\[
 \pi\circ\sigma=\operatorname{id}_{\mathbb Z_5}.
\]

\begin{lemma}\label{lem:maps}
The maps \(\pi:X\to C_5\) and \(\sigma:C_5\to X\) are graph
homomorphisms.
\end{lemma}

\begin{proof}
If \(x\sim_X y\), then
\[
 \pi(y)-\pi(x)=\pi(y-x)\in\{1,-1\},
\]
so \(\pi(x)\sim_{C_5}\pi(y)\).

Conversely, if \(i\sim_{C_5}j\), then \(j-i=\pm1\) in \(\mathbb Z_5\).
Therefore
\[
 \sigma(j)-\sigma(i)=\pm6\in S,
\]
and hence \(\sigma(i)\sim_X\sigma(j)\).
\end{proof}

The image of \(\sigma\), in cyclic order, is
\[
 H=\{0,6,12,3,9\},\qquad 0-6-12-3-9-0.
\]

\begin{lemma}\label{lem:induced}
The subgraph of \(X\) induced by \(H\) is a five-cycle, and
\[
 R=\sigma\circ\pi:X\longrightarrow H
\]
is a retraction.
\end{lemma}

\begin{proof}
For distinct \(i,j\in\mathbb Z_5\),
\[
 \sigma(j)-\sigma(i)=6(j-i)\in S
 \quad\Longleftrightarrow\quad
 j-i\equiv\pm1\pmod5.
\]
Thus the only edges among the five image vertices are the five consecutive
cycle edges, so \(H\) is induced.  By \Cref{lem:maps}, the composite \(R\)
is a graph homomorphism.  For \(h=\sigma(i)\in H\),
\[
 R(h)=\sigma(\pi(\sigma(i)))=\sigma(i)=h,
\]
so \(R\) fixes \(H\) pointwise.
\end{proof}

\section{The five-cycle is a core}

A finite graph is a \emph{core} if every endomorphism is an automorphism.

\begin{lemma}\label{lem:c5core}
Every endomorphism of \(C_5\) is an automorphism.
\end{lemma}

\begin{proof}
Label the vertices by \(\mathbb Z_5\), and let
\(f:C_5\to C_5\) be an endomorphism.  Since consecutive vertices must map
to adjacent vertices,
\[
 f(i+1)-f(i)=\epsilon_i,\qquad
 \epsilon_i\in\{1,-1\}.
\]
Summing over the five edges gives
\[
 \sum_{i=0}^{4}\epsilon_i\equiv0\pmod5.
\]
The sum of five signs is one of
\[
 -5,-3,-1,1,3,5.
\]
Only \(\pm5\) is divisible by \(5\), so all signs are equal.  Consequently
\(f(i)=a+i\) for all \(i\), or \(f(i)=a-i\) for all \(i\).  These are the
rotations and reflections of \(C_5\), and each is an automorphism.
\end{proof}

\begin{proposition}\label{prop:core}
The core of \(X\) is isomorphic to \(C_5\).
\end{proposition}

\begin{proof}
By \Cref{lem:induced}, \(X\) retracts onto \(H\cong C_5\).  Hence \(X\) and
\(C_5\) are homomorphically equivalent.  By \Cref{lem:c5core}, \(C_5\) is
already a core.  The core is unique up to isomorphism within a finite
homomorphic-equivalence class.
\end{proof}

\section{The power-cardinality obstruction}

\begin{theorem}\label{thm:main}
There is a Cayley graph on \(\mathbb Z_{15}\) whose core is not a Cayley
graph on any power of \(\mathbb Z_{15}\).
\end{theorem}

\begin{proof}
Take the graph \(X\) above.  Its core has five vertices by
\Cref{prop:core}.  A Cayley graph on \((\mathbb Z_{15})^m\) has exactly
\(15^m\) vertices.  If \(m=0\) is allowed, this number is \(1\); if
\(m\geq1\), it is divisible by \(3\).  In neither case is it \(5\).
\end{proof}

\begin{remark}[Scope]\label{rem:scope}
The example uses the composite-order cyclic group \(\mathbb Z_{15}\).
It refutes an unrestricted formulation quantifying over arbitrary abelian
groups.  It does not refute variants restricted to
\(\mathbb Z_p\), elementary abelian groups, or cubelike graphs.
\end{remark}

\appendix

\section{Checklist for verification}

\begin{enumerate}
  \item Check that \(S=-S\), \(0\notin S\), and \(1\in S\).
  \item Verify that reduction modulo \(5\) preserves every edge.
  \item Verify that \(i\mapsto6i\) preserves every edge of \(C_5\).
  \item Check that the image \(\{0,6,12,3,9\}\) is induced and that the
  composite fixes it pointwise.
  \item Verify the sign-sum proof that every endomorphism of \(C_5\) is an
  automorphism.
  \item Compare \(5\) with \(15^m\) for every \(m\geq0\).
  \item Keep the arbitrary-abelian statement separate from the prime and
  elementary-abelian variants.
\end{enumerate}

\begin{thebibliography}{9}

\bibitem{HellNesetril}
P.~Hell and J.~Ne\v{s}et\v{r}il,
\emph{Graphs and Homomorphisms},
Oxford Lecture Series in Mathematics and its Applications, vol.~28,
Oxford University Press, 2004.

\bibitem{Samal}
R.~\v{S}\'amal,
\emph{Cores of Cayley graphs},
Open Problem Garden, Problem~125 (2008), with a counterexample observation
by G.~Royle and subsequent scope clarification by the author.

\end{thebibliography}

\end{document}
